\section{Continuing Professional Development (CPD) Log}
\label{app:cpd}

This log records self-directed learning and skills development activities undertaken during the project. Activities are grouped by domain and listed in approximate chronological order.

\subsection*{Domain 1: LoRaWAN Protocol and RF Propagation}

\begin{table}[H]
\centering
\caption{CPD activities: LoRaWAN and RF propagation.}
\label{tab:cpd_lorawan}
\begin{tabularx}{\textwidth}{p{3.5cm} p{2cm} X}
\toprule
\textbf{Activity} & \textbf{Period} & \textbf{Learning Outcome} \\
\midrule
Study of LoRaWAN 1.0.4 specification (LoRa Alliance) & S1 W2--W3 & Understood ADR algorithm, duty cycle enforcement, and the MAC-layer confirmation mechanism; informed correct EMA-ADR implementation. \\
\midrule
Reading: Zorbas et al.\ (2025), ``Revisiting LoRaWAN scalability'' & S1 W3 & Identified the SF Monoculture failure mode as the primary motivation; clarified the distinction between intra-SF and inter-SF collisions. \\
\midrule
Study of Two-Ray Ground Reflection model (Rappaport, \textit{Wireless Communications}) & S1 W4 & Derived the breakpoint distance formula and implemented the piecewise path loss model; verified against published numerical examples. \\
\midrule
Reading: Nakamura et al.\ (2022), LoRa UAV gateway measurements & S1 W4 & Calibrated log-normal shadowing standard deviation ($\sigma = 4$\,dB) against field measurement data; confirmed Two-Ray model suitability for UAV-to-ground links. \\
\midrule
Study of LoRa Capture Effect (Croce et al.) & S1 W6 & Understood the 6\,dB SIR threshold for successful frame capture; implemented Capture Effect resolver as a post-collision arbitration step within the Gymnasium \texttt{step()} function. \\
\bottomrule
\end{tabularx}
\end{table}

\subsection*{Domain 2: Deep Reinforcement Learning and Implementation}

\begin{table}[H]
\centering
\caption{CPD activities: Deep Reinforcement Learning.}
\label{tab:cpd_drl}
\begin{tabularx}{\textwidth}{p{3.5cm} p{2cm} X}
\toprule
\textbf{Activity} & \textbf{Period} & \textbf{Learning Outcome} \\
\midrule
Study of Sutton \& Barto, \textit{Reinforcement Learning: An Introduction} (Ch.\ 6, 9--10) & S1 W1--W2 & Established theoretical foundations: temporal-difference learning, Q-learning convergence, function approximation with neural networks. \\
\midrule
Study of DQN paper (Mnih et al., 2015) and Double DQN (van Hasselt et al., 2016) & S1 W2 & Understood experience replay, target network stabilisation, and the overestimation bias corrected by Double DQN; selected Double DQN for implementation. \\
\midrule
Study of Stable-Baselines3 DQN source code (GitHub) & S1 W5 & Identified target network update schedule, replay buffer implementation, and exploration schedule; enabled correct hyperparameter configuration for curriculum training. \\
\midrule
Reading: Domain Randomisation literature (Tobin et al., 2017; Andrychowicz et al., 2020) & S1 W10--W11 & Understood the theoretical basis for randomising simulation parameters to reduce Sim-to-Real gap; applied to sensor position and count randomisation during training. \\
\midrule
Reading: Curriculum Learning survey (Narvekar et al., 2020) & S1 W11 & Designed the three-stage curriculum (grid sizes 100--500--1000; sensor counts 10--20--40) based on task difficulty ordering principles. \\
\midrule
Study of Welch's $t$-test and effect size analysis & S2 W6 & Applied two-sample unequal-variance $t$-test for significance testing across five random seeds; understood conditions under which the test is applicable. \\
\bottomrule
\end{tabularx}
\end{table}

\subsection*{Domain 3: Academic Writing and Professional Skills}

\begin{table}[H]
\centering
\caption{CPD activities: academic writing and professional skills.}
\label{tab:cpd_writing}
\begin{tabularx}{\textwidth}{p{3.5cm} p{2cm} X}
\toprule
\textbf{Activity} & \textbf{Period} & \textbf{Learning Outcome} \\
\midrule
Attendance: Lecture 2 --- Introduction to LaTex and mathematical writing & S1 W2 & Set up LuaLaTeX compilation pipeline; learnt correct mathematical notation for MDP and reward function formulations. \\
\midrule
Attendance: Lecture 3 --- Writing the Draft Introduction and Methodology & S1 W7 & Structured the introduction to follow the ``funnel'' pattern (broad context $\to$ specific gap $\to$ objectives); applied to the Draft Introduction submission. \\
\midrule
Attendance: Lecture 4 --- Writing the Draft/Final Report & S2 W1 & Learnt how to frame quantitative results with appropriate hedging language and statistical context; applied to the Results chapter. \\
\midrule
Completion of Academic Malpractice Awareness training (CANVAS) & S1 W12 & Completed Short Course, Driving Test, Self Declaration Form, and Tutorial Questions as required by Handbook Section~4. \\
\midrule
Version control with Git (branching, commit discipline) & S1--S2 & Maintained a version-controlled repository throughout the project; separate branches for environment development, training experiments, and evaluation. \\
\bottomrule
\end{tabularx}
\end{table}
